% 第7章：前沿方法与系统集成
\chapter{前沿方法与系统集成}

\begin{levelbox}
本章介绍任务规划的前沿方法（大语言模型、神经符号方法）和系统集成实践，通过综合案例展示完整的任务规划系统设计。建议学时：6学时。
\end{levelbox}

% =============================================
\section{大语言模型与规划}
% =============================================

\subsection{LLM的规划能力与局限}

大语言模型（LLM）在自然语言理解和生成方面表现出色，近年来也被用于任务规划。

\begin{keypoint}[LLM规划的优势与局限]
\textbf{优势}：
\begin{itemize}
  \item 自然语言接口，易于交互
  \item 丰富的世界知识
  \item 能够理解复杂的任务描述
  \item 具有一定的常识推理能力
\end{itemize}

\textbf{局限}：
\begin{itemize}
  \item 规划正确性难以保证（可能产生不可执行的计划）
  \item 复杂约束处理能力有限
  \item 长程规划容易出错
  \item 缺乏形式化验证能力
\end{itemize}
\end{keypoint}

\subsection{LLM+P：结合LLM与规划器}

\begin{definition}[LLM+P框架]
LLM+P利用LLM的语言理解能力和传统规划器的规划能力：
\begin{enumerate}
  \item LLM将自然语言任务描述转换为PDDL问题
  \item 调用Fast Downward等规划器求解
  \item LLM将规划结果转换回自然语言
\end{enumerate}
\end{definition}

\begin{lstlisting}[caption=LLM+P流程示例]
输入：自然语言任务描述
"卫星需要先观测北京区域，然后观测上海区域，最后返回待机状态"

LLM转换为PDDL：
(:goal (and
  (observed beijing)
  (observed shanghai)
  (standby satellite1)
  (before (observed beijing) (observed shanghai))
))

规划器求解，返回动作序列：
1. move(satellite1, pos1, beijing_orbit)
2. observe(satellite1, beijing)
3. move(satellite1, beijing_orbit, shanghai_orbit)
4. observe(satellite1, shanghai)
5. enter_standby(satellite1)

LLM转换回自然语言：
"首先将卫星移动到北京上空轨道位置，执行北京区域观测，
然后转移到上海上空，完成上海区域观测，最后进入待机模式。"
\end{lstlisting}

\subsection{思维链与推理增强}

\textbf{思维链}（Chain-of-Thought, CoT）：让LLM逐步推理，显式展示思考过程。

\textbf{思维树}（Tree of Thoughts, ToT）：扩展CoT，允许分支探索和回溯。

\textbf{ReAct}：结合推理（Reasoning）和行动（Acting），LLM可以调用外部工具获取信息。

\begin{example}[ReAct在任务规划中的应用]
\begin{verbatim}
思考：需要规划卫星观测任务，首先查询目标区域的云量
动作：调用气象API获取北京区域未来24小时云量预报
观测：北京区域预计14:00-16:00云量最低(20%)
思考：应该安排在14:00-16:00窗口观测北京
动作：检查卫星在该时段是否经过北京上空
观测：卫星在14:30-15:00经过北京上空
思考：时间窗口匹配，可以安排观测
动作：生成观测指令...
\end{verbatim}
\end{example}

\subsection{国内研究进展}

\begin{appbox}[title={\textbf{国内LLM规划研究}}]
\begin{itemize}
  \item \textbf{AutoPlan}（秦龙等, 自动化学报2024）：面向复杂任务的自主规划框架，结合LLM与规划引擎
  \item \textbf{MetaGPT}（深度求索DeepWisdom团队）：多智能体软件开发框架
  \item \textbf{ChatDev}（清华NLP实验室）：基于对话的软件开发智能体
  \item 国内大模型（文心一言、通义千问、ChatGLM、DeepSeek）在规划任务上持续优化
\end{itemize}
\end{appbox}

% =============================================
\section{神经符号规划与图神经网络}
% =============================================

\subsection{神经符号方法}

\begin{definition}[神经符号AI]
神经符号AI结合神经网络的学习能力和符号AI的推理能力：
\begin{itemize}
  \item 神经网络：感知、特征提取、模式识别
  \item 符号系统：逻辑推理、规划、可解释性
\end{itemize}
\end{definition}

\textbf{在规划中的应用}：
\begin{itemize}[leftmargin=2em]
  \item 神经网络学习启发式函数
  \item 神经网络生成动作候选
  \item 符号规划器保证解的正确性
\end{itemize}

\subsection{图神经网络在规划中的应用}

规划问题天然具有图结构（状态转移图、任务依赖图），图神经网络（GNN）可以有效捕获这些结构信息。

\textbf{应用场景}：
\begin{itemize}[leftmargin=2em]
  \item 学习PDDL问题的表示
  \item 预测动作的效用
  \item 生成规划启发式
\end{itemize}

% =============================================
\section{【综合案例】完整任务规划系统设计}
% =============================================

\begin{caseboxA}[title={\textbf{综合案例：卫星观测任务规划系统}}]
\textbf{系统目标}：设计一个完整的卫星观测任务规划系统，集成前6章的方法。

\textbf{系统需求}：
\begin{enumerate}
  \item 支持自然语言任务输入
  \item 处理多星协同调度
  \item 考虑不确定性（云层）
  \item 支持对抗性场景分析
  \item 提供可解释的规划结果
\end{enumerate}
\end{caseboxA}

\subsection{系统架构设计}

\begin{figure}[htbp]
\centering
\begin{tikzpicture}[
  box/.style={rectangle, draw, rounded corners, minimum width=3cm, minimum height=0.8cm, align=center},
  arrow/.style={->, thick}
]
  % 输入层
  \node[box, fill=blue!10] (input) at (0,4) {自然语言输入\\LLM解析};

  % 建模层
  \node[box, fill=green!10] (model) at (0,2.5) {PDDL建模\\约束生成};

  % 规划层
  \node[box, fill=yellow!10] (plan1) at (-3,1) {经典规划器\\(Fast Downward)};
  \node[box, fill=yellow!10] (plan2) at (0,1) {HTN规划器\\(SHOP2)};
  \node[box, fill=yellow!10] (plan3) at (3,1) {约束求解器\\(SMT)};

  % 优化层
  \node[box, fill=orange!10] (opt) at (0,-0.5) {多智能体协调\\博弈分析};

  % 输出层
  \node[box, fill=red!10] (output) at (0,-2) {规划结果\\可视化与解释};

  % 箭头
  \draw[arrow] (input) -- (model);
  \draw[arrow] (model) -- (plan1);
  \draw[arrow] (model) -- (plan2);
  \draw[arrow] (model) -- (plan3);
  \draw[arrow] (plan1) -- (opt);
  \draw[arrow] (plan2) -- (opt);
  \draw[arrow] (plan3) -- (opt);
  \draw[arrow] (opt) -- (output);
\end{tikzpicture}
\caption{卫星观测任务规划系统架构}
\end{figure}

\subsection{各模块说明}

\textbf{1. 输入解析模块}
\begin{itemize}[leftmargin=2em]
  \item 使用LLM理解自然语言任务描述
  \item 提取目标、约束、优先级等信息
  \item 生成结构化的任务规格
\end{itemize}

\textbf{2. 问题建模模块}
\begin{itemize}[leftmargin=2em]
  \item 根据任务规格生成PDDL/HDDL模型
  \item 编码时间窗口、资源约束
  \item 考虑不确定性因素的概率模型
\end{itemize}

\textbf{3. 规划求解模块}
\begin{itemize}[leftmargin=2em]
  \item 根据问题特性选择合适的规划器
  \item 简单问题：经典规划器
  \item 层次任务：HTN规划器
  \item 复杂约束：SMT求解器
\end{itemize}

\textbf{4. 多智能体协调模块}
\begin{itemize}[leftmargin=2em]
  \item 协调多星任务分配
  \item 解决资源冲突
  \item 对抗场景下的博弈分析
\end{itemize}

\textbf{5. 输出与解释模块}
\begin{itemize}[leftmargin=2em]
  \item 生成可执行的任务指令
  \item 提供规划结果的可视化
  \item 解释规划决策的理由
\end{itemize}

% =============================================
\section{可解释性与人机协同}
% =============================================

\subsection{可解释规划}

规划结果的可解释性对用户信任和人机协同至关重要。

\textbf{解释类型}：
\begin{enumerate}[leftmargin=2em]
  \item \textbf{为什么}：为什么选择这个动作？
  \item \textbf{为什么不}：为什么不选择另一个动作？
  \item \textbf{假设}：如果条件变化，规划会如何改变？
\end{enumerate}

\subsection{人机协同规划}

\begin{keypoint}[人机协同模式]
\begin{itemize}
  \item \textbf{人在环路}（Human-in-the-Loop）：人类审核和批准规划结果
  \item \textbf{混合主动性}：人和机器根据情况动态分配决策权
  \item \textbf{可调自主性}：根据任务关键程度调整自动化水平
\end{itemize}
\end{keypoint}

% =============================================
\section{规划系统的安全性}
% =============================================

\subsection{安全关键规划}

在安全关键应用中（航空、军事、医疗），规划系统必须满足严格的安全要求。

\textbf{安全性保障方法}：
\begin{enumerate}[leftmargin=2em]
  \item \textbf{形式化验证}：使用模型检测验证规划的正确性
  \item \textbf{安全约束}：硬编码安全边界，不允许违反
  \item \textbf{监控与回退}：运行时监控，异常时回退到安全状态
  \item \textbf{冗余设计}：多个独立规划器交叉验证
\end{enumerate}

\subsection{对抗鲁棒性}

在对抗环境中，需要考虑攻击者可能利用规划系统的漏洞。

\textbf{潜在攻击}：
\begin{itemize}[leftmargin=2em]
  \item 输入扰动：修改任务描述误导规划器
  \item 模型污染：污染学习数据影响学习型规划器
  \item 对抗样本：针对神经网络组件的对抗攻击
\end{itemize}

% =============================================
\section{本章小结与展望}
% =============================================

\subsection*{本章小结}

\begin{itemize}[leftmargin=2em]
  \item LLM为任务规划提供了自然语言接口
  \item LLM+P等混合方法结合了LLM和传统规划器的优势
  \item 神经符号方法是未来规划研究的重要方向
  \item 完整的规划系统需要集成多种方法
  \item 可解释性和安全性是实际部署的关键考量
\end{itemize}

\subsection*{未来展望}

\begin{enumerate}[leftmargin=2em]
  \item \textbf{更强的LLM规划能力}：随着LLM的发展，直接规划能力将增强
  \item \textbf{更紧密的神经-符号融合}：端到端可微的规划系统
  \item \textbf{更广泛的多模态规划}：结合视觉、语言、物理的规划
  \item \textbf{更可靠的安全保障}：形式化验证与AI的结合
\end{enumerate}

\begin{discussbox}
\textbf{讨论题}：
\begin{enumerate}
  \item LLM直接生成规划与LLM+P方法各有什么优缺点？未来会如何发展？
  \item 如果你要设计一个实际的卫星观测规划系统，会如何权衡自动化程度和人工干预？
  \item 规划系统的可解释性和性能之间是否存在权衡？如何平衡？
  \item 任务规划技术在军事应用中应该遵循什么伦理准则？
\end{enumerate}

\textbf{课程项目}：
\begin{enumerate}
  \item 选择案例A或案例B，设计并实现一个简化版的任务规划系统
  \item 系统应至少包含：问题建模、规划求解、结果展示三个模块
  \item 撰写项目报告，说明设计思路、实现方法和测试结果
\end{enumerate}
\end{discussbox}
